\documentclass[11pt]{article}
\usepackage[margin=1in]{geometry}
\usepackage[T1]{fontenc}
\usepackage{lmodern}
\usepackage{amsmath,amssymb,amsthm}
\usepackage{graphicx}
\usepackage[hidelinks]{hyperref}

\newtheorem{theorem}{Theorem}
\newtheorem{lemma}{Lemma}
\newcommand{\F}{\mathbb F}
\newcommand{\R}{\mathbb R}
\newcommand{\C}{\mathbb C}
\newcommand{\Sym}{\operatorname{Sym}}
\newcommand{\ran}{\operatorname{ran}}
\newcommand{\tr}{\operatorname{tr}}
\newcommand{\spn}{\operatorname{span}}

\title{A Three-Plane Counterexample to Problem 5.4\\
of arXiv:1508.06687}
\author{Run fa\_banach\_001, agent\_lane\_05}
\date{11 August 2026}

\begin{document}
\maketitle

\section*{Problem and answer}

Cahill, Casazza, Jasper, and Woodland ask the following converse on page~18,
Problem~5.4, of \emph{Phase Retrieval} \cite{CCJW}.  If a family
$\{W_i\}_{i=1}^M$ of subspaces of $\mathcal H_N$ still fails phase retrieval
after adjoining \emph{every} possible subspace $W_{M+1}$, must there be
orthonormal bases of the $W_i$ whose union has a partition $F_1,F_2$ with
\[
 \dim\spn F_1\le M-2,
 \qquad
 \dim\spn F_2\le M-1?
\]
The answer is no, already for three proper planes in dimension three.  Since
our example has $M=N=3$, it also answers the version in which the two printed
bounds are read as $N-2,N-1$.

\begin{theorem}\label{thm:main}
Let $\F\in\{\R,\C\}$, let $e_1,e_2,e_3$ be the standard basis of $\F^3$,
and set
\[
 W_1=e_1^\perp,
 \qquad W_2=e_2^\perp,
 \qquad W_3=(e_1+e_2+e_3)^\perp.
\]
Then:
\begin{enumerate}
\item for every subspace $W_4\subseteq\F^3$, the four subspaces
      $W_1,W_2,W_3,W_4$ fail phase retrieval;
\item for every choice of orthonormal bases $B_i$ of $W_i$, no partition
      $B_1\mathbin{\dot\cup}B_2\mathbin{\dot\cup}B_3=F_1\mathbin{\dot\cup}F_2$
      satisfies
      \[
        \dim\spn F_1\le1,
        \qquad
        \dim\spn F_2\le2.
      \]
\end{enumerate}
Consequently these three planes are a counterexample to Problem~5.4.
\end{theorem}

\section*{Proof intuition}

Lift a signal $x$ to the rank-one matrix $xx^*$.  Squared projection norms
become linear measurements,
\[
 \|P_Wx\|^2=\tr(P_Wxx^*).
\]
Four measurements cannot avoid the determinant-zero cone inside the
six-dimensional real space of real symmetric $3\times3$ matrices: their
kernel has dimension at least two, and an odd cubic vanishes on every real
two-plane.  Our first two planes rule out semidefinite kernel points, so the
resulting singular point is an indefinite rank-two difference $xx^*-yy^*$.

The partition obstruction has a different geometry.  If its two spans are a
line $L$ and a plane $H$, then each two-dimensional $W_i$ either equals $H$
or has an orthonormal basis split between $L$ and $H$.  The three displayed
planes force one $W_i$ to equal $H$, after which their three pairwise
intersection lines violate the required orthogonality.

\section*{No fourth subspace can work}

Let $P_i$ be the orthogonal projection onto $W_i$, including an arbitrary
$W_4$.  We work in the six-dimensional real vector space $\Sym_3(\R)$ even
when $\F=\C$.  Each Hermitian $P_i$ defines a real linear functional there,
and we set
\[
 \mathcal A(Q)=\bigl(\tr(P_1Q),\ldots,\tr(P_4Q)\bigr).
\]
Thus $\dim_\R\ker\mathcal A\ge2$.  Choose a real two-dimensional subspace
$E\subseteq\ker\mathcal A$.  On the unit circle of $E$, the continuous
function $Q\mapsto\det Q$ is odd:
\[
 \det(-Q)=-\det Q.
\]
The intermediate value theorem therefore gives a nonzero $Q\in E$ with
$\det Q=0$.

There is no nonzero positive semidefinite matrix in $\ker\mathcal A$.  Indeed,
if $Q\ge0$ and $\tr(P_iQ)=0$, then
\[
 0=\tr(P_iQ)=\|P_iQ^{1/2}\|_{\mathrm{HS}}^2,
 \qquad
 \ran Q\subseteq\ker P_i=W_i^\perp.
\]
Using only $i=1,2$ gives
\[
 \ran Q\subseteq W_1^\perp\cap W_2^\perp
 =\spn(e_1)\cap\spn(e_2)=\{0\}.
\]
The same argument applied to $-Q$ excludes nonzero negative semidefinite
matrices.  Hence our nonzero singular $Q$ is indefinite.  It has rank two and
a spectral decomposition
\[
 Q=xx^*-yy^*
\]
for nonzero orthogonal real vectors $x,y$.  Since $Q\in\ker\mathcal A$,
\[
 \|P_ix\|^2-\|P_iy\|^2
 =\tr(P_iQ)=0,
 \qquad i=1,2,3,4.
\]
But orthogonal nonzero $x,y$ cannot differ by a real sign or a complex phase.
Thus the four subspaces fail phase retrieval.  This proof permits an arbitrary
real or complex fourth subspace.

\section*{The requested partition cannot exist}

Assume to the contrary that orthonormal bases $B_i$ and a partition
$F_1,F_2$ exist.  Write
\[
 L=\spn F_1,
 \qquad H=\spn F_2,
 \qquad \dim L\le1,
 \qquad \dim H\le2.
\]
Each $W_i$ is two-dimensional.  Its two basis vectors cannot both lie in
$L$.  If both lie in $H$, then $W_i=H$.  Otherwise exactly one lies in each
part; in that case $L\subset W_i$, the line $H\cap W_i$ is generated by the
$F_2$ basis vector, and
\begin{equation}\label{eq:orthogonality}
 L\perp(H\cap W_i).
\end{equation}

The three normals $e_1,e_2,e_1+e_2+e_3$ are linearly independent, hence
$W_1\cap W_2\cap W_3=\{0\}$.  Therefore not all three bases can split between
$L$ and $H$; at least one, say that of $W_k$, lies wholly in $H$, forcing
$H=W_k$.  The other two planes, say $W_i,W_j$, must split, and their common
$F_1$ line is necessarily
\[
 L=W_i\cap W_j.
\]
There are three possibilities for $k$, all impossible:
\begin{align*}
 k=1:&\quad W_2\cap W_3=\spn(e_1-e_3),
       &W_2\cap W_1=\spn(e_3),
       &(e_1-e_3)^*e_3=-1;\\
 k=2:&\quad W_1\cap W_3=\spn(e_2-e_3),
       &W_1\cap W_2=\spn(e_3),
       &(e_2-e_3)^*e_3=-1;\\
 k=3:&\quad W_1\cap W_2=\spn(e_3),
       &W_1\cap W_3=\spn(e_2-e_3),
       &e_3^*(e_2-e_3)=-1.
\end{align*}
In each row the last nonzero inner product contradicts
\eqref{eq:orthogonality} for the indicated plane.  This proves the second
claim and completes the proof of Theorem~\ref{thm:main}.

\section*{Verification and scope}

The script \texttt{code/verify\_counterexample.py} checks the three projection
matrices and the three intersection-line obstructions exactly.  It also tests
102 rationally generated fourth subspaces of all ranks and finds an
indefinite singular lifted-kernel witness in every case.  This computation is
only a sanity check; the determinant and incidence arguments above prove the
universal statements.

The result is a negative answer, not a classification of one-subspace
completability.  A bounded search for the exact Problem~5.4 and Corollary~5.3
wording found the source but no later resolution; novelty should still receive
expert literature review.

\begin{thebibliography}{9}
\bibitem{CCJW}
J.~Cahill, P.~G. Casazza, J.~Jasper, and L.~M. Woodland,
\emph{Phase Retrieval},
arXiv:1508.06687 (2015),
\href{https://arxiv.org/abs/1508.06687}{arXiv:1508.06687}.
\end{thebibliography}

\end{document}
